```latex
\chapter{Vehicular Networks}
\label{sec:chapter9}

Vehicular communication networks constitute an important component of
Intelligent Transportation Systems (ITS), enabling vehicles to exchange
information with other vehicles, roadside infrastructure, pedestrians, and
communication networks. The integration of wireless communication with
on-board sensing and computing allows vehicles to obtain information beyond
the sensing range of their local sensors and supports applications ranging
from collision avoidance and traffic management to cooperative driving and
passenger infotainment.

Vehicular communication technology has evolved considerably over the past
two decades. Early research primarily focused on Vehicular Ad Hoc Networks
(VANETs), in which vehicles dynamically establish wireless links with nearby
vehicles and roadside infrastructure. Dedicated Short-Range Communications
(DSRC), based primarily on IEEE 802.11p and the Wireless Access in Vehicular
Environments (WAVE) architecture, emerged as one of the earliest standardized
solutions for vehicular communication
\cite{kenney2011dsrc,jiang2008ieee,morgan2010dsrc}.

The rapid development of cellular communication subsequently resulted in
Cellular Vehicle-to-Everything (C-V2X). LTE-V2X introduced direct sidelink
communication between vehicles through the PC5 interface, while the evolution
toward fifth-generation cellular systems resulted in New Radio V2X
(NR-V2X). NR-V2X extends the capabilities of LTE-V2X to support advanced
applications including vehicle platooning, cooperative perception, advanced
driving, and remote driving
\cite{seo2016lte,molinamasegosa2017ltev,lien2020nr}.

This chapter presents the evolution of vehicular networking technologies,
beginning with VANETs and IEEE 802.11p-based communication and subsequently
discussing LTE-V2X and NR-V2X. Particular emphasis is placed on safety
message dissemination, mobility-related challenges, sidelink communication,
and radio resource allocation.

% ============================================================
\section{Vehicular Communication Networks}
\label{sec:vehicular_communication_networks}
% ============================================================

A vehicular network consists of vehicles equipped with communication,
sensing, positioning, and computing capabilities that exchange information
with other entities in the transportation environment. Unlike conventional
wireless networks, vehicular networks are characterized by high node
mobility, rapidly changing network topology, intermittent connectivity,
stringent communication latency requirements, and highly dynamic wireless
channels.

Vehicular communications are commonly grouped under the general term
Vehicle-to-Everything (V2X). Depending on the communicating entities, V2X
communication can be categorized as follows:

\begin{itemize}
    \item \textbf{Vehicle-to-Vehicle (V2V):} Direct communication among
    vehicles for exchanging information such as position, velocity,
    acceleration, heading, and safety warnings.

    \item \textbf{Vehicle-to-Infrastructure (V2I):} Communication between
    vehicles and roadside infrastructure such as Roadside Units (RSUs),
    traffic lights, or traffic-management systems.

    \item \textbf{Vehicle-to-Pedestrian (V2P):} Communication between
    vehicles and vulnerable road users equipped with compatible communication
    devices.

    \item \textbf{Vehicle-to-Network (V2N):} Communication between vehicles
    and remote network services through cellular infrastructure.
\end{itemize}

Together, these communication modes provide vehicles with information beyond
their direct sensing range and enable cooperative decision making among road
users.

% ------------------------------------------------------------
\subsection{Evolution from VANETs to V2X Networks}
% ------------------------------------------------------------

Vehicular Ad Hoc Networks evolved from Mobile Ad Hoc Networks (MANETs), but
the mobility characteristics of vehicles distinguish VANETs from conventional
ad hoc networks. Vehicle movement is constrained by road geometry and traffic
rules, while relative velocities can be very high. Consequently, communication
links may exist only for short durations and the network topology can change
rapidly.

Early VANET architectures relied mainly on distributed communication among
vehicles and RSUs. IEEE 802.11p was developed specifically to support
low-latency wireless access in vehicular environments
\cite{jiang2008ieee,kenney2011dsrc}. The resulting DSRC/WAVE architecture
enabled direct communication without requiring vehicles to first establish a
connection with cellular infrastructure.

The emergence of LTE and device-to-device communication introduced a second
major technology family for V2X. Cellular V2X combines the wide-area
connectivity of cellular networks with direct sidelink communication between
vehicles. The resulting evolution can broadly be represented as

\[
\text{VANET}
\rightarrow
\text{DSRC/IEEE 802.11p}
\rightarrow
\text{LTE-V2X}
\rightarrow
\text{NR-V2X}
\rightarrow
\text{Internet of Vehicles}.
\]

Future vehicular networks increasingly integrate communication, sensing,
computing, cloud services, and edge intelligence, leading toward the broader
concept of the Internet of Vehicles (IoV)
\cite{zhou2020evolutionary,xu2018iov}.

% ------------------------------------------------------------
\subsection{Vehicular Applications}
% ------------------------------------------------------------

Applications supported by vehicular networks can broadly be divided into
safety-related applications, traffic-efficiency applications, and
infotainment services.

Safety applications are generally the most time-critical. Examples include
forward-collision warning, emergency braking notification, intersection
collision warning, lane-change assistance, and hazardous-road-condition
notification. These applications typically require frequent dissemination of
vehicle status information together with highly reliable reception and low
communication latency.

Traffic-efficiency applications use information collected from vehicles and
road infrastructure to improve traffic flow. Examples include congestion
notification, adaptive route planning, traffic signal coordination, and
optimal speed advisory.

Infotainment applications include Internet connectivity, multimedia
streaming, map updates, passenger services, and cloud-based vehicle services.
Unlike many safety applications, these services can require relatively high
data rates while tolerating larger communication delays.

The emergence of automated driving has introduced another category of
communication-intensive applications. Vehicle platooning, cooperative
perception, sensor-information sharing, trajectory coordination, and remote
driving require substantially stronger combinations of reliability, latency,
and throughput than traditional vehicular safety messaging.

% ============================================================
\section{IEEE 802.11-Based Vehicular Networks}
\label{sec:80211_vanet}
% ============================================================

Dedicated Short-Range Communications was among the earliest widely studied
radio technologies for V2X communication. In the United States, DSRC was
developed around operation in the 5.9~GHz band and was supported by the
IEEE 802.11p physical and Medium Access Control (MAC) layers together with
the IEEE 1609 WAVE family of standards
\cite{kenney2011dsrc,morgan2010dsrc}.

IEEE 802.11p evolved from IEEE 802.11a but introduced modifications to enable
communication in rapidly varying vehicular environments
\cite{jiang2008ieee}. Vehicles equipped with On-Board Units (OBUs) can
communicate directly with other vehicles or with RSUs without requiring a
conventional Wi-Fi association procedure.

% ------------------------------------------------------------
\subsection{Safety Message Dissemination}
% ------------------------------------------------------------

A fundamental operation in VANETs is the dissemination of periodic safety
messages. Such messages convey information including vehicle position,
speed, direction, acceleration, and vehicle state. Neighboring vehicles use
these messages to maintain awareness of the surrounding traffic environment.

In DSRC-based networks, safety packets are commonly transmitted through
broadcast communication. This approach minimizes connection-establishment
overhead and allows a single transmission to inform multiple neighboring
vehicles. Vehicle-to-vehicle safety messaging using DSRC has therefore been
an important research topic since the early development of VANETs
\cite{xu2004safety,jiang2006dsrc}.

However, broadcast communication does not inherently guarantee successful
packet reception. As vehicular density increases, a large number of vehicles
may simultaneously compete for channel access. Packet collisions and channel
congestion can consequently reduce the probability that safety information is
successfully delivered.

% ------------------------------------------------------------
\subsection{MAC Challenges in IEEE 802.11p}
% ------------------------------------------------------------

IEEE 802.11p employs contention-based channel access derived from the IEEE
802.11 family. Although this distributed operation is attractive because it
does not require centralized coordination, its performance can degrade in
dense traffic conditions.

Important MAC-layer challenges include:

\begin{itemize}
    \item contention among a large number of simultaneously transmitting
    vehicles;
    \item hidden-terminal collisions;
    \item rapidly varying numbers of neighboring vehicles;
    \item excessive channel occupancy caused by frequent safety beacons;
    \item lack of deterministic packet-access delay; and
    \item degradation in packet-delivery reliability under high network load.
\end{itemize}

These limitations have motivated numerous adaptive MAC, congestion-control,
and mobility-aware access schemes. For example, mobility-aware MAC protocols
have been developed specifically to reduce packet collisions for vehicular
safety applications \cite{lyu2018momac}, while TDMA-based approaches such as
VeMAC have also been proposed to provide more predictable broadcast access
\cite{omar2013vemac}.

% ============================================================
\section{Cellular Vehicle-to-Everything Communications}
\label{sec:cv2x}
% ============================================================

The limitations of contention-based vehicular communication and the rapid
evolution of mobile cellular systems motivated the development of Cellular
Vehicle-to-Everything communication. Cellular systems provide extensive
network coverage, centralized radio-resource management, mature security
mechanisms, and integration with wide-area network services
\cite{araniti2013lte,gyawali2021challenges}.

A major development that enabled cellular V2X was the introduction of direct
device-to-device sidelink communication. C-V2X therefore supports both
infrastructure-based communication through the conventional cellular
interface and direct communication between nearby road users.

% ------------------------------------------------------------
\subsection{LTE-V2X Architecture}
% ------------------------------------------------------------

LTE-V2X was standardized by the 3rd Generation Partnership Project (3GPP)
for supporting vehicular communication. Two major communication interfaces
are relevant:

\begin{itemize}
    \item \textbf{Uu interface:} Conventional communication between a vehicle
    User Equipment (UE) and the cellular base station.

    \item \textbf{PC5 interface:} Direct sidelink communication between
    nearby UEs without routing each packet through the cellular base station.
\end{itemize}

The PC5 sidelink is particularly important for safety applications because
vehicles can directly exchange messages even when conventional cellular
coverage is unavailable.

% ------------------------------------------------------------
\subsection{LTE-V2X Resource Allocation}
% ------------------------------------------------------------

LTE-V2X introduced two principal sidelink resource-allocation mechanisms,
commonly referred to as Mode~3 and Mode~4.

\subsubsection{Mode 3: Network-Scheduled Resource Allocation}

In Mode~3, the serving evolved Node B (eNB) coordinates sidelink radio
resources. The network can therefore use information available at the base
station to reduce conflicting resource assignments.

Two general scheduling strategies can be employed. In dynamic scheduling,
vehicles request radio resources whenever packets need to be transmitted.
Although flexible, frequent requests increase signaling overhead and access
latency.

Semi-Persistent Scheduling (SPS) is more appropriate when messages are
generated periodically. Resources are reserved for several consecutive
transmissions, thereby reducing repeated scheduling requests and control
signaling.

\subsubsection{Mode 4: Autonomous Resource Selection}

Mode~4 allows vehicles to select resources autonomously from a preconfigured
resource pool and can therefore operate outside cellular coverage.

Resource selection employs a sensing-based SPS procedure. Vehicles monitor
sidelink resources and attempt to exclude resources that are likely to be
occupied by neighboring transmitters before selecting candidate transmission
resources. Once selected, the resource can be retained for multiple
transmissions before a resource-reselection procedure is performed.

The autonomous nature of Mode~4 makes it suitable for distributed V2V
communication. However, because vehicles only possess partial information
about the transmission decisions of other UEs, resource collisions remain an
important performance limitation
\cite{gonzalezmartin2018mode4,bazzi2018mode4,molinamasegosa2017mode4}.

% ============================================================
\section{5G and NR-V2X Communications}
\label{sec:nrv2x}
% ============================================================

The requirements of advanced automated-driving applications exceed the
capabilities originally considered during the design of LTE-V2X. NR-V2X was
therefore standardized in 3GPP Release~16 to complement LTE-V2X and support
more advanced vehicular services
\cite{3gpp38885,lien2020nr}.

Examples include vehicle platooning, cooperative and extended sensing,
advanced driving, and remote driving. Compared with LTE-V2X, these
applications can involve a combination of periodic and aperiodic traffic,
variable packet sizes, stringent latency constraints, and substantially larger
information payloads.

NR-V2X also expands the supported communication types. Whereas basic
LTE-V2X sidelink operation primarily focuses on broadcast communication,
NR-V2X supports broadcast, groupcast, and unicast communication. Groupcast is
particularly relevant to applications such as platooning, where information is
intended for a specific group of cooperating vehicles.

% ------------------------------------------------------------
\subsection{Flexible NR Numerology}
% ------------------------------------------------------------

Another important feature of NR is flexible numerology. LTE-V2X employs a
15~kHz subcarrier spacing, whereas NR-V2X can support larger subcarrier
spacings including 30, 60, and 120~kHz. Increasing the subcarrier spacing
reduces the corresponding slot duration and can therefore help support
latency-sensitive applications.

The appropriate numerology nevertheless depends on factors such as vehicle
speed, propagation conditions, delay requirements, and sensitivity to
inter-carrier and inter-symbol interference.

% ------------------------------------------------------------
\subsection{NR-V2X Sidelink Resource Allocation}
% ------------------------------------------------------------

NR-V2X defines two principal resource-allocation modes.

\subsubsection{Mode 1}

In Mode~1, radio resources are assigned through the cellular network. A gNB
can dynamically schedule sidelink resources or provide configured grants for
repeated transmissions. Network coordination permits the infrastructure to
consider communication requirements such as packet periodicity, packet size,
latency, and reliability when assigning resources.

\subsubsection{Mode 2}

Mode~2 supports autonomous resource selection and can operate without
continuous assistance from the gNB. Vehicles select resources from configured
or preconfigured sidelink resource pools.

Unlike the LTE-V2X Mode~4 design, which was principally developed around
periodic traffic, NR-V2X Mode~2 must support both periodic and aperiodic
traffic. Consequently, enhancements to sensing and resource-selection
procedures are important aspects of NR-V2X resource allocation
\cite{ali2021mode2}.

% ============================================================
\section{Mobility Challenges and Safety Message Dissemination}
\label{sec:mobility_challenges}
% ============================================================

Vehicular networking differs fundamentally from many conventional wireless
systems because communication conditions vary rapidly in both space and
time. Several challenges directly influence the reliability of safety-message
dissemination.

\subsection{Rapidly Changing Topology}

Vehicles frequently enter and leave each other's communication range. The
duration of a V2V link may therefore be short, particularly when vehicles
travel in opposite directions or at significantly different speeds.

\subsection{Doppler and Time-Varying Channels}

High vehicle velocities result in significant Doppler shifts and reduced
channel coherence times. Reliable channel estimation and synchronization
become increasingly difficult as relative velocity and carrier frequency
increase.

\subsection{Congestion and Interference}

Dense vehicular scenarios can contain hundreds of simultaneously active
transmitters. In contention-based networks, this produces channel congestion,
whereas sidelink cellular systems must address interference caused by reuse of
the same time-frequency resources.

\subsection{Latency and Reliability}

Safety applications impose substantially more stringent communication
requirements than conventional best-effort wireless services. A packet that
arrives after its safety information has become obsolete can have little
practical value even when it is successfully decoded. Consequently,
vehicular communication protocols must jointly consider packet reception
reliability, latency, packet age, and information-update frequency.

% ============================================================
\section{Radio Resource Allocation in Vehicular Networks}
\label{sec:v2x_resource_allocation}
% ============================================================

Efficient radio resource allocation is particularly important in vehicular
networks because multiple vehicles must frequently share limited
time-frequency resources while satisfying heterogeneous communication
requirements.

Resource-allocation methods can broadly be categorized as centralized,
distributed, and hybrid approaches.

In centralized approaches, a base station or roadside controller collects
network information and determines resource assignments. Such methods can
achieve coordinated interference management but introduce signaling overhead
and depend on infrastructure availability.

Distributed methods allow individual vehicles to select their own
transmission resources. They offer greater robustness to infrastructure
unavailability but must make decisions using incomplete knowledge of
neighboring resource selections.

Spatial resource reuse is particularly important in V2X networks because the
same resource can potentially be assigned to geographically separated
transmitters whose mutual interference is sufficiently small. Consequently,
location-aware, clustering-based, graph-based, and interference-aware
resource-allocation mechanisms have been extensively investigated
\cite{abantoleon2017graph,cecchini2018reuse}.

More recently, learning-based approaches have also been investigated for
vehicular resource allocation. Reinforcement learning enables vehicles or
infrastructure controllers to adapt resource-selection decisions using
observations of network conditions rather than relying entirely on fixed
heuristic rules.

% ============================================================
\section{Heterogeneous Vehicular Networks}
\label{sec:heterogeneous_v2x}
% ============================================================

IEEE 802.11-based V2X and cellular V2X possess complementary characteristics.
DSRC provides distributed short-range communication without requiring
cellular infrastructure, whereas cellular technologies provide extensive
network coverage and coordinated resource management.

This observation motivated heterogeneous architectures in which multiple
radio access technologies coexist. Such systems can dynamically select
between IEEE 802.11p and cellular communication according to application
requirements, channel availability, network load, or vehicle location
\cite{abboud2016interworking,ucar2016hybrid}.

However, heterogeneous V2X introduces additional challenges, including
network selection, vertical handover, mobility management, duplicated
signaling, and coordination among different radio interfaces.

% ============================================================
\section{Evolution Toward the Internet of Vehicles}
\label{sec:iov}
% ============================================================

The Internet of Vehicles represents a broader evolution of conventional
vehicular communication. Instead of treating V2X solely as a mechanism for
short-range safety-message exchange, IoV integrates vehicles with road
infrastructure, cellular networks, edge-computing platforms, cloud systems,
and intelligent transportation services.

Future vehicular networks are expected to combine communication with
distributed sensing, computing, artificial intelligence, and automated
driving. Multi-access edge computing can reduce the latency associated with
remote processing, while network slicing and software-defined networking can
provide differentiated services for heterogeneous vehicular applications.

The resulting architecture moves from simple connectivity toward
cooperative intelligence, where vehicles exchange not only status messages
but also sensor observations, environmental information, trajectories, and
control-related information.


% ============================================================
\section{Research Contributions in Vehicular Networks}
\label{sec:vehicular_research_contributions}
% ============================================================

The research on vehicular networks initially focused on the
performance of IEEE 802.11p-based Vehicular Ad Hoc Networks
(VANETs) for road-safety applications. Subsequent work extended
this study toward cellular V2X technologies, including LTE-V2X
and NR-V2X, with particular emphasis on radio resource
allocation and support for advanced vehicular communication
services.

% ------------------------------------------------------------
\subsection{Traffic-Safety Evaluation of IEEE 802.11p/WAVE Networks}
\label{subsec:wave_safety_evaluation}
% ------------------------------------------------------------

An early study investigated the effectiveness of the Wireless
Access in Vehicular Environments (WAVE) protocol in supporting
traffic-safety applications in Vehicular Ad Hoc Networks
(VANETs) \cite{venkatesh2019wave}. The work considered the
IEEE 802.11p physical layer together with the IEEE 1609 WAVE
protocol stack for communication between vehicles and Roadside
Units (RSUs).

In the considered architecture, vehicles are equipped with
On-Board Units (OBUs) and periodically broadcast beacon
messages containing information such as their position and
intended route. RSUs collect information from vehicles within
their communication range and estimate traffic congestion along
the corresponding road segments. When congestion is detected,
the RSU broadcasts the relevant information to approaching
vehicles, enabling them to select an alternative route.

The study was aimed at understanding how communication-network
parameters affect the ultimate traffic-safety performance of a
vehicular network. In particular, the effects of transmission
bit rate, beacon transmission interval, transmit power, vehicle
density, and vehicle mobility were investigated. Unlike studies
that characterize a VANET exclusively through communication
metrics such as throughput, packet loss, or packet delay, the
work considered the vehicle collision rate as an application-level
measure of traffic safety.

A coupled traffic and communication simulation framework was
used for the evaluation. The Simulation of Urban Mobility
(SUMO) was employed to generate vehicular mobility and road
traffic, while OMNeT++ together with the Veins framework was
used to simulate IEEE 802.11p/WAVE communication. Vehicles
periodically transmit their location information, and RSUs
maintain information about traffic conditions in their coverage
regions. When congestion is identified, vehicles can determine
an alternate path using shortest-path routing.

The evaluation demonstrated the strong coupling between
communication parameters and traffic-safety performance.
Frequent beacon transmission improves the freshness of traffic
information available to vehicles; however, excessively frequent
transmissions can increase MAC-layer contention and packet
collisions. Similarly, increasing transmission power initially
improves communication coverage, but excessive transmit power
increases the number of simultaneously contending vehicles and
can consequently increase MAC-layer collisions.

Vehicle density and mobility were also shown to significantly
affect safety performance. Higher vehicular density increases
both road congestion and contention for the shared wireless
channel. The study therefore highlighted that the communication
parameters of a WAVE-based VANET cannot be configured
independently of the underlying traffic conditions.

This work established the initial connection between
communication-level configuration and road-safety performance
in the vehicular networking studies considered in this chapter.
It also motivated subsequent investigations of more coordinated
cellular V2X technologies and adaptive radio-resource
management mechanisms.

% ------------------------------------------------------------
\subsection{Performance Analysis and Adaptive Resource Allocation in Cellular V2X}
\label{subsec:cv2x_performance_resource_allocation}
% ------------------------------------------------------------

Following the earlier investigations of IEEE 802.11p/WAVE-based
vehicular networks, subsequent research considered the evolution
toward cellular Vehicle-to-Everything (C-V2X) communication
\cite{bhargavi2022rat}. In particular, the performance of
LTE-V2X and New Radio V2X (NR-V2X) was investigated with
emphasis on the impact of physical- and MAC-layer parameters
on the reliability and timeliness of vehicular information
exchange.

LTE-V2X, introduced by 3GPP for direct sidelink vehicular
communication, enables vehicles to communicate through the PC5
interface. In LTE-V2X, a fixed subcarrier spacing of $15$~kHz
is employed, while radio resources can either be assigned by the
cellular infrastructure or selected autonomously by the vehicles.
NR-V2X extends this framework by introducing a flexible OFDM
numerology and additional mechanisms for supporting more
stringent and heterogeneous vehicular communication
requirements.

A detailed simulation-based evaluation was carried out using
LTEV2Vsim for LTE-V2X and WiLabV2Xsim for NR-V2X. The
study investigated the effects of several important system
parameters, including inter-vehicle distance, vehicle density,
Modulation and Coding Scheme (MCS), transmit power, and
Subcarrier Spacing (SCS). Performance was evaluated using
three principal metrics: Packet Reception Ratio (PRR), Data
Age, and Update Delay.

The Packet Reception Ratio characterizes communication
reliability by measuring the fraction of intended receivers that
successfully receive a transmitted packet. Update Delay
characterizes the time separation between consecutive
successfully received beacon messages from a transmitting
vehicle, while Data Age captures the freshness of the
successfully delivered vehicular information. The simultaneous
consideration of reliability and information timeliness is
particularly important for safety-related vehicular applications,
where successful reception alone is insufficient if the received
information has become outdated.

The numerical study showed that the performance of NR-V2X
depends strongly on the selected physical-layer and
resource-allocation parameters. Increasing the distance between
communicating vehicles degrades packet-reception performance,
while vehicle density influences the interference and contention
experienced within the shared sidelink resource pool. Similarly,
the MCS and transmit power determine a trade-off between
communication robustness, spectral efficiency, and interference
to neighboring vehicles.

An important feature of NR-V2X is its flexible numerology.
Unlike LTE-V2X, which employs a fixed $15$~kHz subcarrier
spacing, NR-V2X supports multiple subcarrier spacings. Larger
SCS values shorten the corresponding slot duration and improve
robustness against frequency-domain impairments associated with
high mobility. However, a single SCS configuration is not
necessarily suitable for all vehicular conditions.

Motivated by this observation, an adaptive resource-allocation
mechanism was proposed in which the SCS is dynamically
selected according to the operating conditions of the vehicular
network. Rather than retaining a fixed numerology configuration,
the proposed approach adapts the radio-resource configuration
with the objective of improving the Packet Reception Ratio.
Simulation results demonstrated that the adaptive strategy can
provide improved reception performance compared with the
default resource-allocation configuration.

This study therefore extended the vehicular networking research
from the evaluation of contention-based IEEE 802.11p/WAVE
systems toward adaptive radio-resource management in cellular
V2X. It also highlighted the importance of jointly considering
communication reliability and information freshness when
evaluating resource-allocation strategies for highly dynamic
vehicular networks.

% ------------------------------------------------------------
\subsection{Learning-Based Resource Allocation for Aperiodic NR-V2X Traffic}
\label{subsec:drl_nrv2x_resource_allocation}
% ------------------------------------------------------------

The increasing diversity of V2X applications introduces traffic
patterns that differ considerably from the periodic safety-message
generation originally considered in conventional sidelink scheduling.
In particular, NR-V2X must accommodate aperiodic packet generation
and variable packet sizes while maintaining high reliability and
information freshness. To address these requirements, a learning-based
adaptive resource-allocation framework was developed for NR-V2X
Mode~2 \cite{indulekha2024drl}.

NR-V2X Mode~2 enables vehicles to autonomously select sidelink
resources using Sensing-Based Semi-Persistent Scheduling (SB-SPS).
Under SB-SPS, a vehicle reserves resources for a number of
consecutive transmissions according to a Resource Reservation
Interval (RRI). When the reselection counter expires, the vehicle
either retains the previously selected resource with probability
$P_{\mathrm{keep}}$ or performs a new resource-selection procedure.
Candidate resources are identified through sidelink sensing using
Sidelink Control Information (SCI), Reference Signal Received Power
(RSRP), and Received Signal Strength Indicator (RSSI).

Although SB-SPS is effective for periodic message generation, its
semi-persistent reservation mechanism can lead to inefficient resource
utilization when traffic becomes aperiodic. If the next packet is
generated later than the previously reserved transmission opportunity,
the reserved resource may remain unused. Similarly, variations in
packet size can cause a previously reserved resource to become
insufficient for a subsequent larger packet, forcing the vehicle to
perform resource reselection.

To improve the adaptability of sidelink resource reservation, three
important parameters were jointly considered: the Modulation and
Coding Scheme (MCS), the Resource Reservation Interval (RRI), and
the resource Keep-Alive probability $P_{\mathrm{keep}}$. The MCS
determines the trade-off between transmission rate and link
robustness, while the RRI controls the temporal spacing between
reserved transmission opportunities. The parameter
$P_{\mathrm{keep}}$ determines whether previously selected sidelink
resources are retained after expiration of the reselection counter.

A Double Deep Q-Network (DDQN)-based resource-allocation
mechanism was developed to dynamically adapt these parameters
according to the observed vehicular-network conditions. The
optimization objective balances communication reliability,
information freshness, and channel congestion through a utility of
the form
\begin{equation}
    \Phi =
    \gamma_{1}\,\mathrm{PRR}
    -
    \gamma_{2}\,\mathrm{DA}
    -
    \gamma_{3}\,\mathrm{CBR},
\end{equation}
where $\mathrm{PRR}$ denotes the Packet Reception Ratio,
$\mathrm{DA}$ denotes the Data Age, $\mathrm{CBR}$ denotes the
Channel Busy Ratio, and $\gamma_{1}$, $\gamma_{2}$, and
$\gamma_{3}$ determine their relative importance.

The corresponding reinforcement-learning problem is represented as
a Markov Decision Process. The state contains information describing
the current channel occupancy, packet-transmission status, timing
information, and vehicle-position updates. The action taken by the
learning agent consists of selecting the triplet
\begin{equation}
    a_t =
    \left[
    \mathrm{MCS}(t),
    \mathrm{RRI}(t),
    P_{\mathrm{keep}}(t)
    \right].
\end{equation}
The reward is derived from the reliability--freshness--congestion
utility $\Phi$, enabling the agent to learn resource configurations
that respond to changes in packet generation, packet size, and
network conditions.

DDQN was adopted to improve the stability of value-based learning
by separating action selection from target-value evaluation. A
behavior Q-network determines the resource-allocation action, while
a target network is periodically updated to provide stable learning
targets. Experience replay is used to train the network from
previously observed state--action--reward transitions.

The proposed framework was evaluated using WiLabV2Xsim under
NR-V2X Mode~2 operation. The experiments considered variations in
vehicle density, packet size, packet-generation interval, MCS,
RRI, and resource keep probability. The results showed that
dynamically adapting these parameters provides a better balance
between packet-reception reliability and information freshness than
using fixed parameter configurations.

In particular, adaptive MCS selection improved packet reception at
larger communication distances while reducing Data Age for
different packet sizes. Adaptive selection of $P_{\mathrm{keep}}$
improved robustness under changing vehicle densities, while dynamic
RRI adaptation was particularly beneficial for aperiodic traffic,
where a fixed reservation interval can result in either unused
resources or delayed information updates.

This work extended the earlier parameter-based resource-allocation
studies toward intelligent and environment-aware radio-resource
management. It demonstrated that learning-based adaptation can be
used to modify fundamental NR-V2X Mode~2 scheduling parameters in
response to heterogeneous and unpredictable vehicular traffic.

% ------------------------------------------------------------
\subsection{Learning-Based Resource Allocation for Aperiodic NR-V2X Traffic}
\label{subsec:drl_nrv2x_resource_allocation}
% ------------------------------------------------------------

The increasing diversity of V2X applications introduces traffic
patterns that differ considerably from the periodic safety-message
generation originally considered in conventional sidelink scheduling.
In particular, NR-V2X must accommodate aperiodic packet generation
and variable packet sizes while maintaining high reliability and
information freshness. To address these requirements, a learning-based
adaptive resource-allocation framework was developed for NR-V2X
Mode~2 \cite{indulekha2024drl}.

NR-V2X Mode~2 enables vehicles to autonomously select sidelink
resources using Sensing-Based Semi-Persistent Scheduling (SB-SPS).
Under SB-SPS, a vehicle reserves resources for a number of
consecutive transmissions according to a Resource Reservation
Interval (RRI). When the reselection counter expires, the vehicle
either retains the previously selected resource with probability
$P_{\mathrm{keep}}$ or performs a new resource-selection procedure.
Candidate resources are identified through sidelink sensing using
Sidelink Control Information (SCI), Reference Signal Received Power
(RSRP), and Received Signal Strength Indicator (RSSI).

Although SB-SPS is effective for periodic message generation, its
semi-persistent reservation mechanism can lead to inefficient resource
utilization when traffic becomes aperiodic. If the next packet is
generated later than the previously reserved transmission opportunity,
the reserved resource may remain unused. Similarly, variations in
packet size can cause a previously reserved resource to become
insufficient for a subsequent larger packet, forcing the vehicle to
perform resource reselection.

To improve the adaptability of sidelink resource reservation, three
important parameters were jointly considered: the Modulation and
Coding Scheme (MCS), the Resource Reservation Interval (RRI), and
the resource Keep-Alive probability $P_{\mathrm{keep}}$. The MCS
determines the trade-off between transmission rate and link
robustness, while the RRI controls the temporal spacing between
reserved transmission opportunities. The parameter
$P_{\mathrm{keep}}$ determines whether previously selected sidelink
resources are retained after expiration of the reselection counter.

A Double Deep Q-Network (DDQN)-based resource-allocation
mechanism was developed to dynamically adapt these parameters
according to the observed vehicular-network conditions. The
optimization objective balances communication reliability,
information freshness, and channel congestion through a utility of
the form
\begin{equation}
    \Phi =
    \gamma_{1}\,\mathrm{PRR}
    -
    \gamma_{2}\,\mathrm{DA}
    -
    \gamma_{3}\,\mathrm{CBR},
\end{equation}
where $\mathrm{PRR}$ denotes the Packet Reception Ratio,
$\mathrm{DA}$ denotes the Data Age, $\mathrm{CBR}$ denotes the
Channel Busy Ratio, and $\gamma_{1}$, $\gamma_{2}$, and
$\gamma_{3}$ determine their relative importance.

The corresponding reinforcement-learning problem is represented as
a Markov Decision Process. The state contains information describing
the current channel occupancy, packet-transmission status, timing
information, and vehicle-position updates. The action taken by the
learning agent consists of selecting the triplet
\begin{equation}
    a_t =
    \left[
    \mathrm{MCS}(t),
    \mathrm{RRI}(t),
    P_{\mathrm{keep}}(t)
    \right].
\end{equation}
The reward is derived from the reliability--freshness--congestion
utility $\Phi$, enabling the agent to learn resource configurations
that respond to changes in packet generation, packet size, and
network conditions.

DDQN was adopted to improve the stability of value-based learning
by separating action selection from target-value evaluation. A
behavior Q-network determines the resource-allocation action, while
a target network is periodically updated to provide stable learning
targets. Experience replay is used to train the network from
previously observed state--action--reward transitions.

The proposed framework was evaluated using WiLabV2Xsim under
NR-V2X Mode~2 operation. The experiments considered variations in
vehicle density, packet size, packet-generation interval, MCS,
RRI, and resource keep probability. The results showed that
dynamically adapting these parameters provides a better balance
between packet-reception reliability and information freshness than
using fixed parameter configurations.

In particular, adaptive MCS selection improved packet reception at
larger communication distances while reducing Data Age for
different packet sizes. Adaptive selection of $P_{\mathrm{keep}}$
improved robustness under changing vehicle densities, while dynamic
RRI adaptation was particularly beneficial for aperiodic traffic,
where a fixed reservation interval can result in either unused
resources or delayed information updates.

This work extended the earlier parameter-based resource-allocation
studies toward intelligent and environment-aware radio-resource
management. It demonstrated that learning-based adaptation can be
used to modify fundamental NR-V2X Mode~2 scheduling parameters in
response to heterogeneous and unpredictable vehicular traffic.

% ------------------------------------------------------------
\subsection{Dynamic Resource and Power Allocation for ISAC-Aided NR-V2X}
\label{subsec:isac_nrv2x_resource_power}
% ------------------------------------------------------------

As vehicular communication systems evolve toward higher levels
of automation, communication reliability alone is no longer
sufficient to characterize system performance. Advanced V2X
applications increasingly require vehicles to simultaneously
exchange information and sense their surrounding environment.
Motivated by this requirement, an Integrated Sensing and
Communication (ISAC)-aided NR-V2X framework was developed
for jointly optimizing communication and sensing performance
\cite{indulekha2024isac}.

The proposed framework considers an NR-V2X sidelink network
in which vehicles exchange Cooperative Awareness Messages
(CAMs) while simultaneously performing radar sensing using
the reflected communication waveform. This enables the same
radio resources and waveform structure to support both
information exchange and estimation of neighboring vehicle
parameters such as range and relative velocity.

For sensing, an OFDM-based radar model is employed. The
transmitted V2X waveform is reflected by neighboring vehicles,
and the received echo is processed to estimate propagation delay
and Doppler shift. These parameters are subsequently used to
estimate the target range and relative velocity. A Delay--Doppler
representation is generated from the received signal, and target
peaks are detected using a Cell-Averaging Constant False Alarm
Rate (CA-CFAR) procedure.

The communication component follows the NR-V2X sidelink
resource-allocation framework. In Mode~2, vehicles autonomously
select transmission resources using the Sensing-Based
Semi-Persistent Scheduling (SB-SPS) mechanism. Candidate
resources are identified using Sidelink Control Information
(SCI), Reference Signal Received Power (RSRP), and Received
Signal Strength Indicator (RSSI) measurements. The performance
of this resource-allocation process depends strongly on physical-
and MAC-layer parameters such as the Modulation and Coding
Scheme (MCS), Subcarrier Spacing (SCS), and transmission power.

A key challenge in an ISAC-enabled vehicular network is that
the parameters that improve communication performance do not
necessarily provide the best sensing performance. For example,
changes in transmit power and SCS affect not only packet
reception and channel utilization but also the accuracy of range
and velocity estimation. Hence, the problem was formulated as
a joint communication--sensing optimization.

The communication performance is characterized using Packet
Reception Ratio (PRR) and Channel Busy Ratio (CBR), while
the sensing performance is quantified using the Cram\'er--Rao
Lower Bounds (CRLBs) for range and velocity estimation. The
joint utility is expressed as
\begin{equation}
    \Psi =
    \delta_{1}\mathrm{PRR}
    -
    \delta_{2}\mathrm{CBR}
    -
    \delta_{3}
    \left(
    \mathrm{CRLB}_{\mathrm{range}}
    +
    \mathrm{CRLB}_{\mathrm{speed}}
    \right),
\end{equation}
where $\delta_{1}$, $\delta_{2}$, and $\delta_{3}$ control the
relative importance of communication reliability, congestion,
and sensing accuracy.

The corresponding resource- and power-allocation problem is
formulated as
\begin{equation}
    \max_{M_k,\Delta f,P_T} \; \Psi,
\end{equation}
subject to feasible ranges for the MCS $M_k$, subcarrier
spacing $\Delta f$, and transmission power $P_T$.

A Double Deep Q-Network (DDQN) is employed to dynamically
adapt these parameters according to the instantaneous vehicular
environment. The Markov Decision Process state incorporates
packet-transmission information, congestion information, sensing
quality, and timing information. The action at time $t$ is given
by
\begin{equation}
    a_t =
    \left[
    M_k(t),
    \Delta f(t),
    P_T(t)
    \right].
\end{equation}

The reward is directly derived from the joint ISAC utility
$\Psi$, allowing the learning agent to seek operating points
that simultaneously maintain high packet-reception reliability,
low channel congestion, and accurate sensing. Experience replay
and separate primary and target Q-networks are used to improve
learning stability.

The proposed method was evaluated using an extended
WiLabV2Xsim environment incorporating both NR-V2X
communication and sensing functionality. Simulations considered
different vehicle densities, packet sizes, MCS values, subcarrier
spacings, and transmission powers. Results showed that the
learning-based approach was able to outperform fixed
resource-allocation configurations in both communication and
sensing metrics.

In particular, the proposed strategy maintained higher PRR and
lower CBR under increasing vehicular density while
simultaneously reducing the CRLB of range and velocity
estimation. This work therefore extended the earlier
communication-centric resource-allocation framework toward a
joint communication--sensing design, demonstrating the potential
of learning-based ISAC resource management for future
vehicular networks.

% ============================================================
\section{Overview of Contributions}
\label{sec:vehicular_contributions}
% ============================================================

Research carried out in the Mobile Computing Lab has addressed different
aspects of vehicular communication networks, with particular emphasis on
communication under high mobility, safety-message dissemination, and radio
resource management.

\begin{itemize}

    \item \textbf{Dhananjay Kumar:}
    Research on routing protocols and dynamic access mechanisms for vehicular
    networks, with emphasis on communication under rapidly changing network
    topologies.

    \item \textbf{Indulekha:}
    Research on safety-message dissemination and radio resource allocation in
    V2X networks, addressing the reliability and latency requirements of
    vehicular communication.

\end{itemize}

These contributions reflect the broader evolution of vehicular networking
research from distributed VANET communication toward cellular-assisted and
resource-aware V2X systems.

% ============================================================
\section{Summary}
\label{sec:vehicular_summary}
% ============================================================

Vehicular networking has progressed from distributed VANET architectures
toward increasingly sophisticated V2X systems. IEEE 802.11p and DSRC
established the foundation for direct short-range vehicular communication,
particularly for safety-message dissemination. However, contention,
congestion, and limited coordination motivated the investigation of
cellular-based alternatives.

LTE-V2X introduced direct PC5 sidelink communication together with both
network-controlled and autonomous resource allocation. NR-V2X subsequently
extended this framework to support advanced applications, flexible
numerologies, unicast and groupcast communication, and more sophisticated
resource-allocation mechanisms.

The continuing evolution toward the Internet of Vehicles is integrating V2X
communication with sensing, computing, artificial intelligence, and automated
driving. Reliable safety-message dissemination and efficient radio resource
allocation therefore remain central problems in the design of future
vehicular communication networks.

